% ============================================================
% Preamble: paket, header/footer, gaya kode, dan makro bantu
% ============================================================

% ---------- Paket dasar ----------
\usepackage[utf8]{inputenc}
\usepackage[T1]{fontenc}
\usepackage[indonesian]{babel}
\usepackage[a4paper,top=3cm,bottom=3cm,left=3cm,right=3cm]{geometry}

% ---------- Matematika & gambar ----------
\usepackage{amsmath,amssymb}
\usepackage{graphicx}
\usepackage{caption}
\usepackage{subcaption}
\usepackage{float}

% ---------- Tabel ----------
\usepackage{longtable}
\usepackage{array}
\usepackage{booktabs}
\usepackage{tabularx}
\usepackage{multirow}
\usepackage{etoolbox}

\captionsetup{font=footnotesize}
\newif\iftabelkecil
\AtBeginEnvironment{table}{\tabelkeciltrue}
\AtBeginEnvironment{longtable}{\tabelkeciltrue\footnotesize}
\AtBeginEnvironment{tabular}{\iftabelkecil\footnotesize\fi}
\AtBeginEnvironment{tabularx}{\iftabelkecil\footnotesize\fi}
\newcolumntype{Y}{>{\raggedright\arraybackslash}X}
\usepackage[table]{xcolor}

% ---------- Lain-lain ----------
\usepackage[most]{tcolorbox}
\usepackage{forest}
\usepackage{fancyhdr}
\usepackage{titlesec}
\usepackage{listings}
\usepackage{hyperref}


% ---------- Lokasi gambar ----------
% public/figures untuk gambar khusus laporan, ../../output/* untuk hasil run program
\graphicspath{{public/}{public/figures/}{../../output/}{../../output/eda/}{../../output/dtl/}{../../output/svm/}}

% ---------- Gaya source code ----------
\definecolor{codebg}{RGB}{248,248,248}
\definecolor{codecomment}{RGB}{106,115,125}
\definecolor{codekeyword}{RGB}{0,92,197}
\definecolor{codestring}{RGB}{3,47,98}
\lstdefinestyle{py}{
	language=Python,
	backgroundcolor=\color{codebg},
	basicstyle=\ttfamily\footnotesize,
	keywordstyle=\color{codekeyword}\bfseries,
	commentstyle=\color{codecomment}\itshape,
	stringstyle=\color{codestring},
	numbers=left,
	numberstyle=\tiny\color{gray},
	breaklines=true,
	frame=single,
	rulecolor=\color{gray!40},
	showstringspaces=false,
	tabsize=4,
	captionpos=b,
}
\lstset{style=py}
\renewcommand{\lstlistingname}{Kode}
\renewcommand{\lstlistlistingname}{Daftar Kode}

% Sisipkan potongan source code langsung dari repo.
% Pemakaian: \kode{../../src/algoritma/dtl/criteria.py}{6}{14}{Fungsi entropy}{lst:dtl-entropy}
\newcommand{\kode}[5]{\lstinputlisting[firstline=#2,lastline=#3,caption={#4},label={#5}]{#1}}

% ---------- Penanda bagian yang belum ditulis ----------
% Ubah \toggletodotrue menjadi \toggletodofalse untuk menyembunyikan semua TODO saat final.
\newif\iftoggletodo
\toggletodotrue
\newcommand{\todo}[1]{\iftoggletodo\textcolor{red}{\textbf{[TODO:} #1\textbf{]}}\fi}
% Kotak pengganti gambar yang belum ada. Ganti dengan \includegraphics saat gambar tersedia.
\newcommand{\placeholdergambar}[1]{\fbox{\parbox[c][4cm][c]{0.8\textwidth}{\centering\todo{Gambar: #1}}}}

% Subsubbab ikut diberi nomor (mis. 2.3.1.2) dan daftar isi sampai level subbab
\setcounter{secnumdepth}{3}
\setcounter{tocdepth}{2}

% ---------- Identitas dokumen (dipakai di cover & header) ----------
\newcommand{\matkul}{IF3170 Inteligensi Artifisial}
\newcommand{\judullaporan}{Laporan Tugas Besar 2: \textit{Supervised Learning}}
\newcommand{\semester}{Semester I 2026/2027}
\newcommand{\namakelompok}{TreeIsAllYouNeed}
\newcommand{\nomorkelompok}{G04}

% ---------- Header & footer ----------
\setlength{\headheight}{50pt}
\setlength{\headsep}{25pt}
\setlength{\footskip}{28pt}

\newcommand{\kepala}{%
	\fancyhf{}%
	\fancyhead[L]{%
		\includegraphics[height=1.1cm]{logo-itb}\hspace{0.3cm}%
		\includegraphics[height=1.1cm]{logo-stei}%
	}%
	\fancyhead[R]{%
		\begin{minipage}{9cm}
			\raggedleft\bfseries \matkul\\
			\normalsize Laporan Tugas Besar 2\\
			\itshape \semester
		\end{minipage}%
	}%
	\renewcommand{\headrulewidth}{0.4pt}%
	\fancyfoot[L]{\small Kelompok \nomorkelompok}%
	\fancyfoot[C]{\small Halaman \thepage}%
	\renewcommand{\footrulewidth}{0.4pt}%
}
\pagestyle{fancy}
\kepala
\fancypagestyle{plain}{\kepala}

% ---------- Pemenggalan baris ----------
% Nama fungsi/kolom panjang di \texttt boleh dipatah setelah garis bawah, dan spasi antarkata
% boleh melebar sedikit agar tidak ada baris yang keluar margin.
\let\oldunderscore\_
\renewcommand{\_}{\oldunderscore\allowbreak}
\setlength{\emergencystretch}{3em}

% ---------- Hyperref ----------
\hypersetup{colorlinks=true, linkcolor=black, citecolor=blue, urlcolor=blue}
